\section{Product-tripod rigidity and the unavoidable mixing divisor}
\label{pr-section-20260907}
\noindent\textbf{Status.} The geometric statements in this section are ordinary
mathematical theorems, checked by a separate adversarial research agent. The
companion Lean file verifies ten integer algebra statements, including a
two-coordinate matrix separation theorem. It does not encode the geometric
classification or prove the standard $abc$ conjecture.

Let $k$ be a field of characteristic zero, $n\ge1$, and
\[
 U=\mathbb P^1_k\setminus\{0,1,\infty\},\qquad
 R_n=k[x_1,\ldots,x_n,x_1^{-1},(1-x_1)^{-1},\ldots,
 x_n^{-1},(1-x_n)^{-1}].
\]
The one-dimensional Riemann--Hurwitz obstruction does not by itself classify
maps from a fixed product of tripods. The following theorem supplies that
separate higher-dimensional obstruction.

\begin{lemma}[A rectangle and its complement]
\label{pr-rectangle}
In an integral domain, suppose
\[
 AD=BC,\qquad(1-A)(1-D)=(1-B)(1-C).
\]
Then $A+D=B+C$ and $(A-B)(A-C)=0$. Consequently either $A=B$, $C=D$,
or $A=C$, $B=D$.
\end{lemma}
\begin{proof}
Expand the second equality and cancel the first to obtain $A+D=B+C$.
Substitute $D=B+C-A$ in $AD=BC$ and factor. The zero-product law and
the additive equality give the two alternatives.
\end{proof}

\begin{theorem}[Product-tripod rigidity]
\label{pr-rigidity}
Every nonconstant regular $k$-morphism $U^n\to U$ is one of
\[
 x_i,\quad 1-x_i,\quad x_i^{-1},\quad (1-x_i)^{-1},\quad
 \frac{x_i}{x_i-1},\quad\frac{x_i-1}{x_i}
\]
for some $i$. Every dominant regular endomorphism of $U^n$ is therefore
a coordinate permutation followed by these transformations coordinate by
coordinate, and is an automorphism.
\end{theorem}
\begin{proof}
The polynomial ring is a unique factorization domain. After inverting the
irreducibles $x_i,1-x_i$, every unit is
$c\prod_i x_i^{a_i}(1-x_i)^{b_i}$ with $c\in k^*$ and $a_i,b_i\in\mathbb Z$.
Indeed a different irreducible cannot divide a unit or its inverse.
A map to $U$ is exactly a function $f\in R_n$ such that $f$ and $1-f$ are
units. Thus both have separated product expressions
\[
 f=c\prod_i u_i(x_i),\qquad 1-f=d\prod_i v_i(x_i).
\]
Pass to an algebraic closure $K$. If $u_i$ is nonconstant, choose
$s_i,t_i\in U(K)$ with $u_i(s_i)\ne u_i(t_i)$. Fix $j\ne i$, arbitrary
$s_j,t_j\in U(K)$, and all other coordinates. The four values of $f$
at the rectangle, with rows indexed by $s_i,t_i$, satisfy both determinant
identities of Lemma~\ref{pr-rectangle}. Its first-column
entries differ because all factors are nonzero. The lemma therefore forces
each row to be constant. Cancelling the other nonzero factors gives
$u_j(s_j)=u_j(t_j)$. Infinitude of $K$ implies that $u_j$ is constant as
a rational function. Thus a nonconstant $f$ depends on one coordinate only.

Set every other coordinate equal to $2$; this is a valid localization
retraction in characteristic zero. It shows explicitly that $f$ and $1-f$
are units in the corresponding one-variable localization. Write $f=P/Q$
in lowest terms. The nonzero polynomials $P,Q,Q-P$ are pairwise coprime,
and every root of their product lies in $\{0,1\}$. Its reduced radical has
degree at most two. Mason--Stothers yields
$\max(\deg P,\deg Q,\deg(Q-P))<2$; its derivative-zero alternative
would make all three polynomials constant in characteristic zero.
The nonconstant map therefore has degree one and permutes the three
boundary points. Their six permutations give the displayed list.

Finally, dominance of an endomorphism requires its output functions to be
algebraically independent. None can be constant and no two can read the
same input coordinate. The resulting coordinate assignment is a permutation;
the inverse coordinate assignment and inverse boundary permutations are
regular on $U^n$.
\end{proof}

The one-active-coordinate argument is characteristic-free after passage to
an algebraic closure. Degree-one rigidity is not: in characteristic $p$,
$f=x_i^p$ is a regular map to $U$, since $1-f=(1-x_i)^p$, and has degree
$p$ in that coordinate. This is a complete-premise counterexample only to
the characteristic-free strengthening.

\begin{corollary}[Codimension-two deletion does not create an amplifier]
\label{pr-codimension-two}
If $Z\subset U^n$ is closed of codimension at least two, every regular map
$U^n\setminus Z\to U$ extends uniquely to $U^n$ and satisfies
Theorem~\ref{pr-rigidity}.
\end{corollary}
\begin{proof}
The rational functions $f,1/f,1-f,1/(1-f)$ are regular at the generic point
of every irreducible divisor of $\operatorname{Spec}R_n$, since that point
is outside $Z$. In a reduced fraction in the factorial ring $R_n$, this
excludes every nonunit irreducible denominator factor. All four functions
therefore belong to $R_n$. Hence $f$ and $1-f$ are units and define the
extension; uniqueness follows from equality in the function field.
\end{proof}

This closes the specific gate of globally regular, genuinely mixed or
degree-amplifying maps of a fixed finite product of tripods with no new
boundary. It does not close the routes through correspondences, proper
subvarieties, non-product moduli spaces, moving maps, or explicitly charged
new divisors.

\subsection{The exact arithmetic cost of a minimal mixed map}
The function $f(x,y)=xy$ defines a mixed map
$U^2\setminus\{xy=1\}\to U$. At the primitive positive point
$(x,y)=(a/c,b/c)$ with $a+b=c$, put
\[
 N=c^2-ab=a^2+ab+b^2.
\]
The output is the primitive triple $(ab,N,c^2)$. A prime dividing $N$
and one of $a,b,c$ would divide a second entry of the original triple,
by reducing the displayed formula modulo the first. Thus
\[
 \gcd(N,abc)=1,\qquad
 \operatorname{rad}(abNc^2)=\operatorname{rad}(abc)\operatorname{rad}(N),
 \qquad \frac34c^2\le N<c^2.
\]
The size inequalities follow from $(a-b)^2\ge0$ and $ab>0$.
Writing $h=\log c$, $r=\log\operatorname{rad}(abc)$, and
$s=\log\operatorname{rad}(N)$, the output quality and excess are exactly
\[
 q'=\frac{2h}{r+s},\qquad
 E'_\varepsilon
 =2E_\varepsilon+(1+\varepsilon)(r-s),
 \quad E_\varepsilon=h-(1+\varepsilon)r.
\]
The new factor has quadratic size and no overlap with the original radical.
Its radical need not have quadratic size. Any argument using this mixed map
must supply independent arithmetic control of $s$ or another way to pay for
the new divisor. No such uniform control is proved or assumed here.

The formal file \texttt{ProductTripodRectangle.lean} verifies the integer
rectangle identities, the row-or-column dichotomy, a global one-coordinate
dependence theorem for integer-valued matrices with those identities, and
the four displayed mixed-map algebra and size statements. Its ten axiom
queries report only \texttt{propext}, \texttt{Classical.choice}, and
\texttt{Quot.sound}. Unit classification, scheme extension, and the full
geometric theorem remain outside that formal scope.

